\clearpage
\VerdierPhysicalPageAnchor{0144}%
\begingroup
\small
\setlength{\parskip}{0pt}%
\setlength{\emergencystretch}{1.5em}%
\setlength{\abovedisplayskip}{3pt plus 1pt minus 1pt}%
\setlength{\belowdisplayskip}{3pt plus 1pt minus 1pt}%

\VerdierSemanticAnchor{ch2-2-3-1-proof-c-bis-converse}
\noindent 饱和的。由 d) 的证明，子范畴 $Q(Q^{-1}(\mathcal C))$ 严格满且
饱和。而范畴 $Q(Q^{-1}(\mathcal C))$ 是 $\mathcal C$ 的严格满闭包。因此
$\mathcal C$ 是饱和的。这就证明了断言 c)$^{\mathrm{bis}}$。

\VerdierSemanticAnchor{ch2-2-3-1-proof-d-bis}
还需证明最后一个断言。显然，当 $\mathcal C$ 严格满时，
$Q(Q^{-1}(\mathcal C))=\mathcal C$。设 $\mathcal A$ 是包含
$\widehat{\mathcal B}$ 的严格满三角子范畴。需要证明
$\mathcal A=Q^{-1}(Q(\mathcal A))$。把 $\mathcal B$ 替换为
$\widehat{\mathcal B}$ 不会改变范畴 $\mathcal D/\mathcal B$ (2.2.11)，
故可假设 $\mathcal B$ 饱和，从而 $\mathcal B=\widehat{\mathcal B}$。
显然有包含关系 $\mathcal A\subset Q^{-1}(Q(\mathcal A))$，因此只需证明：
$\mathcal D$ 中在 $\mathcal D/\mathcal B$ 里同构于 $\mathcal A$ 的某个对象
的任意对象，仍是 $\mathcal A$ 的对象。

\VerdierSemanticAnchor{ch2-2-3-1-proof-d-bis-roof-setup}
于是设 $X$ 是 $\mathcal A$ 的对象，$Y$ 是 $\mathcal D$ 的对象，并设
$m\colon Q(X)\xrightarrow{\sim}Q(Y)$ 为同构。态射 $m$ 可分解为
$m=Q(f)Q(s)^{-1}$，其中 $s\in S_{\mathcal D}(\mathcal B)$ (2.2.2)，
且 $Q(f)$ 是同构。由于 $\mathcal B$ 饱和，故
$f\in S_{\mathcal D}(\mathcal B)$（(2.1.8) 与 (2.2.4)）。因此在
$\mathcal D$ 中有图式：

\VerdierSemanticAnchor{ch2-2-3-1-proof-d-bis-roof}
\begin{center}
  \begin{tikzcd}[column sep=4.2em,row sep=2.5em,baseline=(current bounding box.center)]
    & Z \arrow[dl,"s"'] \arrow[dr,"f"] & \\
    X && Y
  \end{tikzcd}
  \qquad $s,f\in S_{\mathcal D}(\mathcal B)$\quad .
\end{center}

\VerdierSemanticAnchor{ch2-2-3-1-proof-d-bis-complete}
因 $s$ 属于 $S_{\mathcal D}(\mathcal B)$，存在特异三角
$(Z,X,W,s,u,v)$，其中 $W$ 是 $\mathcal B$ 的对象 (2.1.7)。因此 $W$
是 $\mathcal A$ 的对象；又因 $\mathcal A$ 严格满，$Z$ 也是
$\mathcal A$ 的对象 (1.2.4)。同一论证表明 $Y$ 也是 $\mathcal A$ 的对象，
命题得证。

\VerdierSemanticAnchor{ch2-2-3-2}
\medskip\noindent\textbf{2.3.2.}\quad
\VerdierSemanticAnchor{ch2-2-3-2-free-object}
设 $F\colon C\longrightarrow C'$ 为函子，$X$ 为 $C'$ 的对象。回顾：若存在
态射 $u\colon F(Y)\longrightarrow X$，使得对每个态射
$F(Z)\xrightarrow{v}X$，都存在唯一态射 $w\colon Z\longrightarrow Y$ 使下图
交换，则称 $C$ 的对象 $Y$ \emph{在 $X$ 上 $F$-右自由}，或简称
\emph{在 $X$ 上右自由}：

\VerdierSemanticAnchor{ch2-2-3-2-free-object-diagram}
\begin{center}
  \begin{tikzcd}[column sep=4.7em,row sep=2.7em]
    F(Y) \arrow[rr,"u"] && X \\
    & F(Z) \arrow[ul,"{F(w)}"] \arrow[ur,"v"'] &
  \end{tikzcd}
  \quad .
\end{center}

\VerdierSemanticAnchor{ch2-2-3-2-freedom-arrow-boundary}
箭头 $u\colon F(Y)\longrightarrow X$ 称为二元组 $(X,Y)$ 的\emph{自由箭头}
（相对于函子 $F$）。还可在取一个
\endgroup
